\documentclass[11pt]{article}
\usepackage[a4paper,margin=25mm]{geometry}
\usepackage{algorithm}
\usepackage{algpseudocode}
\usepackage{float}
\usepackage{booktabs}
\usepackage{hyperref}

% \title{Config Generator}
% \author{ROS2-Monitor-Infra Design Document}
% \date{2026-06-15}

\begin{document}
% \maketitle

\section{Config Generator}

The generator creates ROS2-Monitor-Infra runtime YAML from explicit user
choices. The user does not manually write YAML files. Instead, the generator
uses selected ROS2 resources, selected custom converter and verdict classes,
and the selected deployment mode to produce the corresponding runtime YAML.

\begin{enumerate}
  \item discover visible ROS2 topics, services, and actions;
  \item let the user select the resources to monitor;
  \item let the user select custom converter and verdict classes;
  \item combine user choices, deployment template, and selected classes;
  \item write one or more runtime YAML files.
\end{enumerate}

\section{Generation Inputs}

\begin{tabular}{@{}p{0.27\linewidth}p{0.66\linewidth}@{}}
\toprule
Input & Description \\
\midrule
ROS2 resources &
Topics, services, and actions selected by the user from the live ROS2 graph. \\
Converter class &
A \texttt{DataConverter} class selected from the available classes under
\texttt{custom/}. \\
Verdict class &
A \texttt{VerdictService} class selected from the available classes under
\texttt{custom/}. \\
Class parameters &
Optional constructor parameters for the selected custom classes. \\
Deployment mode &
Integrated local, central over MQTT, hybrid, choreographed verdict
aggregation, or offline replay. \\
Transport settings &
MQTT broker and topics, or replay-file path, when required by the chosen
deployment mode. \\
Output settings &
Output directory, session prefix, raw-record output, and verdict exporters. \\
Optional framework settings &
Topic QoS, action phases, and optional built-in transformers.
\\
\bottomrule
\end{tabular}

\clearpage
\section{Config Generation Algorithm}

\begin{algorithm}[H]
\caption{Generate runtime YAML}
\begin{algorithmic}[1]
\Require running ROS2 system and user choices
\Ensure generated configuration files $C$
\State $G \gets \Call{ScanROSResources}{}$
\State \Call{PresentROSResources}{$G$}
\State $S \gets \Call{SelectROSResources}{G}$
\State $P \gets \Call{ScanCustomClasses}{\texttt{custom/}}$
\State $(conv,verdict) \gets \Call{SelectCustomClasses}{P}$
\State $A \gets \Call{CollectClassParameters}{conv,verdict}$
\State $D \gets \Call{SelectDeploymentMode}{}$
\State $O \gets \Call{SelectOutputsAndTransports}{D}$
\State $C \gets \Call{CreateEmptyConfigs}{D,O}$
\ForAll{resource $s$ in $S$}
  \State $entry \gets \Call{MakeSourceEntry}{s.name,s.kind,s.type}$
  \State \Call{ApplyOptionalFrameworkSettings}{entry,s}
  \State \Call{AddSourceToConfigs}{C,entry,D}
\EndFor
\Statex
\State $chain \gets \Call{MakeConverterChain}{conv,verdict,A}$
\State $chain.inputs \gets$ names of selected resources in $S$
\State \Call{AddChainToConfigs}{C,chain,D}
\State \Call{ApplyDeploymentTemplate}{C,D,O}
\State \Call{WriteYAMLFiles}{C}
\State \Return $C$
\end{algorithmic}
\end{algorithm}

\section{Runtime Resource Scanning}

\begin{algorithm}[H]
\caption{Scan ROS2 resources}
\begin{algorithmic}[1]
\Ensure resource catalogue $G$
\State $T \gets$ run \texttt{ros2 topic list -t}
\State $V \gets$ run \texttt{ros2 service list -t}
\State $A \gets$ run \texttt{ros2 action list -t}
\State $G \gets \emptyset$
\ForAll{topic $(name,type)$ in $T$}
  \State add $(topic,name,type)$ to $G$
\EndFor
\ForAll{service $(name,type)$ in $V$}
  \State add $(service,name,type)$ to $G$
\EndFor
\ForAll{action $(name,type)$ in $A$}
  \State add $(action,name,type)$ to $G$
\EndFor
\State \Return sorted $G$
\end{algorithmic}
\end{algorithm}

\clearpage
\section{Configuration Assembly}

The generated YAML is assembled from three parts: selected ROS2 resources,
the selected monitoring classes, and the selected deployment template.

\subsection{ROS2 resource entries}

Each selected ROS2 resource is copied into the corresponding runtime YAML
section:

\begin{itemize}
  \item selected topics become entries under \texttt{topics};
  \item selected services become entries under \texttt{services};
  \item selected actions become entries under \texttt{actions}.
\end{itemize}

\begin{verbatim}
monitor:
  output_dir: <selected output directory>
  session_id_prefix: <selected session prefix>

topics:
  - name: /selected/topic
    type: package/msg/Message
    qos: <optional QoS>

services:
  - name: /selected/service
    type: package/srv/Service

actions:
  - name: /selected/action
    type: package/action/Action
    phases: [feedback, status]
\end{verbatim}

\subsection{Monitoring chain}

The selected converter and verdict classes are written into the
\texttt{converters} section:

\begin{verbatim}
converters:
  - type: custom.my_case.converter:MyConverter
    inputs: [/selected/topic]
    <converter parameters>
    verdict:
      type: custom.my_case.verdict:MyVerdict
      <verdict parameters>
\end{verbatim}

The generator inserts the selected resource names into \texttt{inputs}. The
custom converter decides how to interpret the resulting records.

\subsection{Class parameters}

Constructor parameters selected for the converter are written next to the
converter \texttt{type}. Parameters selected for the verdict service are
written inside the nested \texttt{verdict} block:

\begin{verbatim}
converters:
  - type: custom.rule_based_converter:RuleBasedConverter
    source_match: "^/odom$"
    field_map:
      speed: twist.twist.linear.x
    verdict:
      type: custom.threshold_verdict:ThresholdVerdict
      property_id: speed_limit
      field: speed
      op: ">"
      threshold: 0.3
\end{verbatim}

\subsection{Resource-to-chain binding}

The \texttt{inputs} field binds selected ROS2 resources to the selected
converter. It is a source-name filter, not a semantic description of the
property. For a multi-resource converter, multiple selected resources can be
listed:

\begin{verbatim}
converters:
  - type: custom.fleet_distance:FleetDistanceConverter
    inputs: [/robot1/odom, /robot2/odom]
    robot_a: /robot1
    robot_b: /robot2
    verdict:
      type: custom.fleet_distance:MinimumFleetDistanceVerdict
      minimum_distance: 1.0
\end{verbatim}

\subsection{Deployment wiring}

The deployment template adds the communication and output blocks required by
the selected architecture. For example, a central MQTT deployment adds an MQTT
exporter on the robot side and an MQTT source on the verifier side:

\begin{verbatim}
# robot-side config
exporters:
  - type: mqtt
    broker: <selected broker>
    topic_prefix: <selected topic prefix>

# verifier-side config
verdict_runner:
  source:
    type: mqtt
    broker: <selected broker>
    topic_filter: <selected topic filter>
\end{verbatim}

For local deployment, this wiring can be omitted because the selected sources
and monitoring chain live in the same runtime configuration.

\section{Deployment Output Patterns}

\begin{tabular}{@{}p{0.22\linewidth}p{0.28\linewidth}p{0.42\linewidth}@{}}
\toprule
Mode & Generated files & Main YAML sections \\
\midrule
Integrated local &
\texttt{config.yaml} &
\texttt{monitor}, selected \texttt{topics/services/actions},
\texttt{converters}. \\
Central over MQTT &
\texttt{robot\_config.yaml}, \texttt{verifier\_config.yaml} &
Robot file: selected resources and MQTT \texttt{exporters}.
Verifier file: \texttt{verdict\_runner.source} and \texttt{converters}. \\
Hybrid &
\texttt{robot\_config.yaml}, \texttt{verifier\_config.yaml} &
Robot file: selected resources, local \texttt{converters}, and MQTT
\texttt{exporters}. Verifier file: MQTT source and central
\texttt{converters}. \\
Choreographed &
\texttt{robotN.yaml}, \texttt{aggregator.yaml} &
Robot files: selected local resources and local verdict exporters.
Aggregator file: verdict source and aggregation \texttt{converters}. \\
Offline replay &
\texttt{recorder.yaml}, \texttt{replay.yaml} &
Recorder file: selected resources and file exporter.
Replay file: file source and \texttt{converters}. \\
\bottomrule
\end{tabular}

These patterns are templates. Concrete paths, MQTT topics, and exporter
settings are filled from the selected transport and output settings.

\end{document}
